\subsection*{一步一步建立直观}

\paragraph{1. 一般输入会同时改变长度与方向}

把绿色向量拖到普通方向，红色输出通常不与它共线。若反复作用，方向也会不断调整。寻找特征向量的动机，就是找出能把这种二维混合简化为一维倍率变化的特殊方向。



\paragraph{2. 保持的是直线，不一定是箭头朝向}

当输入落在虚线上，输出与它共线，满足 $Av=\lambda v$。正倍率保留朝向，负倍率反转朝向，零倍率把非零向量压成零；三者都符合特征方程。定义要求 $v$ 非零，否则任意数都会满足方程。



\paragraph{3. 用另一方向补全分解}

右图的实对称例子有两条互相垂直的特征方向，可把任意输入分成这两份。反复作用时，两份坐标分别乘倍率的幂。这个例子的完整分解不能自动推广为“每个实矩阵都有两条实特征方向”，旋转矩阵就是需要注意的情形。



\subsection*{定义、命题与证明}

\paragraph{定义：特征向量与特征值}

若非零向量 $v$ 满足 $Av=\lambda v$，称 $v$ 为 $A$ 的特征向量，$\lambda$ 为对应特征值。允许 $\lambda=0$；禁止把零向量当作特征向量。

\paragraph{命题：一例对称矩阵的两个特征方向}

矩阵 $A=\begin{pmatrix}2&1\\1&2\end{pmatrix}$ 沿 $(1,1)$ 的倍率为三，沿 $(1,-1)$ 的倍率为一。

\begin{proof}

直接计算 $A(1,1)=(3,3)=3(1,1)$，以及 $A(1,-1)=(1,-1)$。两个输入都非零，故满足定义。

它们线性无关，任何平面向量都可分解为这两个方向的组合。在线性作用下，分别把两个分量乘以三和一再相加即可得到输出。这是此矩阵可对角化的具体解释，不代表所有矩阵都具有完整特征基。

\end{proof}
